%% arara directives
% arara: xelatex
% arara: bibtex
% arara: xelatex
% arara: xelatex

%\documentclass{article} % One-column default
\documentclass[twocolumn, switch]{article} % Method A for two-column formatting

% TODO: this package is out of date, before submission I need something else
\usepackage{preprint}
\usepackage{cool}
\usepackage{xparse}

% Save the original cool derivatives if you ever need them
\let\CoolD\D
\let\CoolDt\Dt

\RenewDocumentCommand{\D}{ m g }{%
    \IfNoValueTF{#2}
        {\frac{\partial}{\partial #1}}%
        {\frac{\partial #1}{\partial #2}}%
}

\NewDocumentCommand{\Dt}{ m g }{%
  \IfNoValueTF{#2}
    {\frac{d}{d #1}}%
    {\frac{d #1}{d #2}}%
}


%% Math packages
\usepackage{amsmath, amsthm, amssymb, amsfonts}

%% Bibliography options
\usepackage[numbers,square]{natbib}
\bibliographystyle{unsrtnat}
%\usepackage{natbib}
%\bibliographystyle{Geology}

%% General packages
\usepackage[utf8]{inputenc}	% allow utf-8 input
\usepackage[T1]{fontenc}	% use 8-bit T1 fonts
\usepackage{xcolor}		% colors for hyperlinks
\usepackage[colorlinks = true,
            linkcolor = purple,
            urlcolor  = blue,
            citecolor = cyan,
            anchorcolor = black]{hyperref}	% Color links to references, figures, etc.
\usepackage{booktabs} 		% professional-quality tables
\usepackage{nicefrac}		% compact symbols for 1/2, etc.
\usepackage{microtype}		% microtypography
\usepackage{lineno}		% Line numbers
\usepackage{float}			% Allows for figures within multicol
%\usepackage{multicol}		% Multiple columns (Method B)

\usepackage{lipsum}		%  Filler text

 %% Special figure caption options
\usepackage{newfloat}
\DeclareFloatingEnvironment[name={Supplementary Figure}]{suppfigure}
\usepackage{sidecap}
\sidecaptionvpos{figure}{c}

% Section title spacing  options
\usepackage{titlesec}
\titlespacing\section{0pt}{12pt plus 3pt minus 3pt}{1pt plus 1pt minus 1pt}
\titlespacing\subsection{0pt}{10pt plus 3pt minus 3pt}{1pt plus 1pt minus 1pt}
\titlespacing\subsubsection{0pt}{8pt plus 3pt minus 3pt}{1pt plus 1pt minus 1pt}

% ORCiD insertion
\usepackage{tikz,xcolor,hyperref}

\definecolor{lime}{HTML}{A6CE39}
\DeclareRobustCommand{\orcidicon}{
	\begin{tikzpicture}
	\draw[lime, fill=lime] (0,0) 
	circle [radius=0.16] 
	node[white] {{\fontfamily{qag}\selectfont \tiny ID}};
	\draw[white, fill=white] (-0.0625,0.095) 
	circle [radius=0.007];
	\end{tikzpicture}
	\hspace{-2mm}
}
\foreach \x in {A, ..., Z}{\expandafter\xdef\csname orcid\x\endcsname{\noexpand\href{https://orcid.org/\csname orcidauthor\x\endcsname}
			{\noexpand\orcidicon}}
}
% Define the ORCID iD command for each author separately. Here done for two authors.
\newcommand{\orcidauthorA}{0000-0003-0042-3583}
\newcommand{\orcidauthorB}{0000-0001-6654-6661}
% \newcommand{\orcidauthorC}{0000-0000-0000-0003}
% \newcommand{\orcidauthorD}{0000-0000-0000-0004}

%%%%%%%%%%%%%%%%   Title   %%%%%%%%%%%%%%%%
\title{Acoustic Vector-reflectivity inversion using first-order parametrization}



% Add watermark with submission status
\usepackage{xwatermark}
% Left watermark
\newwatermark[firstpage,color=gray!60,angle=90,scale=0.32, xpos=-4.05in,ypos=0]{\href{https://doi.org/}{\color{gray}{Publication doi}}}
% Right watermark
\newwatermark[firstpage,color=gray!60,angle=90,scale=0.32, xpos=3.9in,ypos=0]{\href{https://doi.org/}{\color{gray}{Preprint doi}}}
% Bottom watermark
\newwatermark[firstpage,color=gray!90,angle=0,scale=0.28, xpos=0in,ypos=-5in]{*correspondence: \texttt{saraivaq@ualberta.ca}}

%%%%%%%%%%%%%%%  Author list  %%%%%%%%%%%%%%%
\usepackage{authblk}
\renewcommand*{\Authfont}{\bfseries}
\author[1]{Átila Saraiva Quintela Soares\orcidA{}}
\author[1]{Mauricio Sacchi\orcidB{}}

\affil[1]{Department of Physics, University of Alberta}
% \affil[2]{Department of Physics, University of Alberta}

% Option 2 for author list
%\author{
%  David S.~Hippocampus\thanks{Use footnote for providing further
%    information about author (webpage, alternative
%    address)---\emph{not} for acknowledging funding agencies.} \\
%  Department of Computer Science\\
%  Cranberry-Lemon University\\
%  Pittsburgh, PA 15213 \\
%  \texttt{hippo@cs.cranberry-lemon.edu} \\
%  %% examples of more authors
%   \And
% Elias D.~Striatum \\
%  Department of Electrical Engineering\\
%  Mount-Sheikh University\\
%  Santa Narimana, Levand \\
%  \texttt{stariate@ee.mount-sheikh.edu} \\
%  \AND
%  Coauthor \\
%  Affiliation \\
%  Address \\
%  \texttt{email} \\
%  % etc.
%}

%%%%%%%%%%%%%%    Front matter    %%%%%%%%%%%%%%
\begin{document}

\twocolumn[ % Method A for two-column formatting
  \begin{@twocolumnfalse} % Method A for two-column formatting
  
\maketitle

\begin{abstract}
% Lorem ipsum dolor sit amet, consectetuer adipiscing elit. Ut purus elit, vestibulum ut, placerat ac, adipiscing vitae, felis. Curabitur dictum gravida mauris. Nam arcu libero, nonummy eget, consectetuer id, vulputate a, magna. Donec vehicula augue eu neque. Pellentesque habitant morbi tristique senectus et netus et malesuada fames ac turpis egestas. Mauris ut leo. Cras viverra metus rhoncus sem. Nulla et lectus vestibulum urna fringilla ultrices. Phasellus eu tellus sit amet tortor gravida placerat. Integer sapien est, iaculis in, pretium quis, viverra ac, nunc. Praesent eget sem vel leo ultrices bibendum. Aenean faucibus. Morbi dolor nulla, malesuada eu, pulvinar at, mollis ac, nulla. Curabitur auctor semper nulla. Donec varius orci eget risus. Duis nibh mi, congue eu, accumsan eleifend, sagittis quis, diam. Duis eget orci sit amet orci dignissim rutrum. Nam dui ligula, fringilla a, euismod sodales, sollicitudin vel, wisi. Morbi auctor lorem non justo. 
\end{abstract}
%\keywords{First keyword \and Second keyword \and More} % (optional)
\vspace{0.35cm}

  \end{@twocolumnfalse} % Method A for two-column formatting
] % Method A for two-column formatting

%\begin{multicols}{2} % Method B for two-column formatting (doesn't play well with line numbers), comment out if using method A


%%%%%%%%%%%%%%%  Main text   %%%%%%%%%%%%%%%
% \linenumbers

\section{Introduction}



\section{Methods}


\section{Results}


\section{Discussion}


\section{Conclusion}


%%%%%%%%%%%% Supplementary Methods %%%%%%%%%%%%
%\footnotesize
%\section*{Methods}

%%%%%%%%%%%%% Acknowledgements %%%%%%%%%%%%%
%\footnotesize
\section*{Acknowledgements}

%%%%%%%%%%%%%%   Bibliography   %%%%%%%%%%%%%%
\normalsize
% TODO: remember to remove URL from the references
\bibliography{references}

\appendix

\section{Derivation of the adjoint-state equations and derivative kernel}

\Style{DDisplayFunc=outset,DShorten=false}

This appendix shows the calculations of the gradient in greater detail using
the adjoint state method based on the works by 
\citet{plessix_review_2006, fichtner_adjoint_2006, cea_conception_1986} .
The objective is to calculate the functional derivative of the following
cost function:
\begin{equation}\label{eq:cost}
    J(c, \mathbf{R}, p_{s}, \mathbf{v}_{s}) = 
    \frac{1}{2 N_{s} N_{r} T} \Sum{ \Sum{
        \Int{
            \left( S_{s,r} p_{s}(\mathbf{x}, t) - d_{s,r}(t) \right)^2
        }{t, 0, T} 
    }{r, 1, N_{r}} }{s, 1, N_{s}},
\end{equation}
where $c(\mathbf{x})$ is the velocity model, 
$\mathbf{R} = \left\{ R_1(\mathbf{x}), R_2(\mathbf{x}) \right\}$ is the
vector-reflectivity model.
% $N_{s}$ is the number of sources, $N_{r}$ is the number of receivers, 
Moreover, $T$ is the recording time, $\mathbf{x} = \left\{ x_1, x_2 \right\} \in \Omega$ are the
spatial coordinates. Additionally, the state variable
$p_{s}(\mathbf{x}, t)$ is the pressure wavefield,
and particle momentum $\mathbf{v}_{s} = \left\{ v_{s, 1}, v_{s, 2} \right\}$, where
$s \in \left\{ 1, \cdots, N_s \right\}$ is the shot index and $r \in \left\{ 1, \cdots, N_{r} \right\}$
is the receiver index. $S_{s,r}$ is a sampling operator that samples
the $p_{s}$ at a specific receiver position $r$ and can be different for every
shot index $s$.
The definition of the sampling operator used in this paper is:
\begin{equation}
    S_{s,r} p_{s} = \Int{p_{s}(\mathbf{x}, t) \, \delta(\mathbf{x} -  \mathbf{x}_{s,r})}{\mathbf{x}, \Omega},
\end{equation}
where $\mathbf{x}_{r}$ is the position of the receiver.

Both velocity ($c$) and vector-reflectivity ($\mathbf{R}$) are related to the state variables 
$p_{s}$ and $v_{s,i}$ by the system of partial diffential equations derived in
\citet{soares_first-order_2025} in form of functional mappings:
\begin{align}\label{eq:pde}
\Style{DDisplayFunc=inset,DShorten=false}
 &   \pderiv[1]{v_{s,i}}{t} + \pderiv[1]{p_{s}}{x_{i}} = 0\\
 &  \pderiv[1]{p_{s}}{t} + \Sum{
        c \pderiv[1]{c}{x_{i}} \, v_{s,i} - 2 c^2 R_{i} + c^2 \pderiv[1]{v_{s, i}}{x_{i}}
    }{i, 1, 2} - f = 0,
\end{align}
where if $F_1$ and $F_2$ are equal to zero, then the $p_s$ and $v_{s,i}$ are
actual physical solutions to this system.
In order to make the solution unique, let us use these boundary and initial
conditions
\begin{align}\label{eq:boundary}
&   p_{s}(\mathbf{x}, 0) = 0,\\
&   v_{s,i}(\mathbf{x}, 0) = 0,\\
&   p_{s}(\mathbf{x}, t) = 0 \quad \text{at  }  \mathbf{x} \in \partial \Omega,\\
&   v_{s,i}(\mathbf{x}, t) = 0 \quad \text{at  }  \mathbf{x} \in \partial \Omega.
\end{align}
All of these equations introduce a new problem, how to calculate the derivative
of the cost function in Eq.~\ref{eq:cost} when we have an implicit relationship
between the state variables and parameters in the form of a PDE? This is solved
by using the adjoint state method.

The first step in the adjoint state method is to assemble the augmented Lagrangian
that turns the system of PDEs and boundary and initial conditions into
regularization terms:
\begin{equation}\label{eq:lagrangian}
    \begin{aligned}
 &   \mathcal{L}\left(c, \mathbf{R}, p_{s}, \mathbf{v}_{s}, \mathbf{\xi}_{s}, \eta_{s}, \gamma_{s}, \mu_{s}, \mathbf{\zeta}_{s}, \mathbf{\lambda_{s}} \right) =
    J(c, \mathbf{R}, p_{s}, \mathbf{v}_{s}) \\
& \quad - 
    \Sum{\Sum{
        \Int{\Int{
            \xi_{s,i} \left( \pderiv[1]{v_{s,i}}{t} + \pderiv[1]{p_{s}}{x_{i}} \right)
        }{t,0,T}}{\mathbf{x}, \Omega}
    }{i,1,2}}{s,1,N_{s}}\\
 & \quad -
    \Sum{
        \Int{\Int{
            \eta_{s} \left( \pderiv[1]{p_{s}}{t} + \Sum{
                            c \pderiv[1]{c}{x_{i}} \, v_{s,i} - 2 c^2 R_{i} + c^2 \pderiv[1]{v_{s, i}}{x_{i}} - f 
                        }{i, 1, N_{s}} \right)
        }{t,0,T}}{\mathbf{x}, \Omega}
    }{s,1,N_{s}}\\
& \quad -
    \Sum{
        \Int{\Int{
            \gamma_{s}(\mathbf{x}, t) p_{s}(\mathbf{x}, t)
        }{t, 0, T}}{\mathbf{x}, \partial\Omega}
    }{s,1,N_{s}}
  -
    \Sum{
        \Int{
            \mu_{s}(\mathbf{x}) p_{s}(\mathbf{x}, 0)
        }{t, 0, T}
    }{s,1,N_{s}}\\
& \quad -
    \Sum{
        \Int{\Int{
            \zeta_{s,i}(\mathbf{x}, t) v_{s,i}(\mathbf{x}, t)
        }{t, 0, T}}{\mathbf{x}, \partial\Omega}
    }{s,1,N_{s}}
  -
    \Sum{
        \Int{
            \lambda_{s,i}(\mathbf{x}) v_{s,i}(\mathbf{x}, 0)
        }{t, 0, T}
    }{s,1,N_{s}}\\
    \end{aligned},
\end{equation}
with $\xi_{s,i}$, $\eta_{s}$ being the adjoint state variables for the PDEs, and
$\gamma_{s}$, $\mu_{s}$, $\zeta_{s,i}$, and $\lambda_{s,i}$ are the adjoint
state variables for the initial and boundary conditions.

The next step is to calculate the directional functional derivative of the
Lagrangian, which we use the definition from \citet[p. 228]{zeidler_applied_1995}:
\begin{equation}
    \CoolD{F}{u}\, h = \lim_{\epsilon \to 0} \frac{F(u + \epsilon h) - F(u)}{\epsilon} = \left[ \frac{d}{d \epsilon} F(u + \epsilon h) \right]_{\epsilon = 0},
\end{equation}
where $\CoolD{F}{u}$ is a linear operator. In order to facilitate the mathematical constraints
and avoid lack of rigor, we will go ahead and assume that $F(u)$ will be always 
continuous, which means that $\CoolD{F}{u}$ will remain the same in all directions
and that it exists. 
Let us also define a ``regular'' and ``partial'' variants
for the multivariate case. With the regular we assume that all the variables
are dependent on each other, which means that the directional functional derivative
becomes:
\begin{multline}\label{eq:regular}
\Style{DDisplayFunc=outset,DShorten=false}
   \CoolD{F(m, u_1, \cdots, u_{N})}{m}\, h = \\
  \quad \left[ \frac{d}{d \epsilon} F(m + \epsilon h, u_1(m + \epsilon h), \cdots , u_N(m + \epsilon h)) \right]_{\epsilon = 0},
\end{multline}
while in the partial case one only considers that specific variable:
\begin{equation}\label{eq:partial}
   \pderiv{F(m, u_1, \cdots, u_{N})}{m}\, h = 
   \left[ \frac{d}{dt} F(m + \epsilon h, u_1, \cdots , u_N) \right]_{\epsilon = 0}.
\end{equation}
These definitions are important because the regular derivative of the cost
function with relation to the model is equal to the partial derivative of the
lagrangian in Eq.~\ref{eq:lagrangian}, so long that the partial derivatives
with relation to all the other variables is set to zero. This can be expressed as
\begin{multline}
    \CoolD{J(c, \mathbf{R}, p_{s}, \mathbf{v}_{s})}{\mathbf{R}} = 
    \pderiv{
\mathcal{L}\left(c, \mathbf{R}, p_{s}, \mathbf{v}_{s}, \mathbf{\xi}_{s}, \eta_{s}, \gamma_{s}, \mu_{s}, \mathbf{\zeta}_{s}, \mathbf{\lambda_{s}} \right)
    }{\mathbf{R}}
    \iff \\
    \pderiv{\mathcal{L}}{\alpha}=0 \quad \text{ for } 
    \alpha \in
\left\{ p_{s}, \mathbf{v}_{s}, \mathbf{\xi}_{s}, \eta_{s}, \gamma_{s}, \mu_{s}, \mathbf{\zeta}_{s}, \mathbf{\lambda_{s}} \right\}.
\end{multline}
In this case, one would also have to set $\CoolD{\mathcal{L}}{c}=0$, but that does
not lead to an equation that needs solving on top of the adjoint state equations
given by \( \pderiv{\mathcal{L}}{\alpha}=0 \),
and this paper assumes that $\mathbf{R}$ and $c$ are not related.

The adjoint state equations are calculated by equaling the derivatives
the augmented Lagrangian (Eq.~\ref{eq:lagrangian}) with relation to the state variables
$p_{s}$ and $v_{s,i}$ to zero. The derivative with relation to $p_{s}$ is calculated as
follows using the definition:
\begin{multline}\label{eq:dLdp}
    \pderiv{\mathcal{L}\left(c, \mathbf{R}, p_{s}, \mathbf{v}_{s}, \mathbf{\xi}_{s}, \eta_{s}, \gamma_{s}, \mu_{s}, \mathbf{\zeta}_{s}, \mathbf{\lambda_{s}} \right)}{p_{s}}\, h_{s} =
    \pderiv{J(c, \mathbf{R}, p_{s}, \mathbf{v}_{s})}{p_{s}}\, h_{s}\\
    - 
    \Sum{\Sum{
        \Int{\Int{
            \xi_{s,i} \left( \pderiv[1]{h_{s}}{x_{i}} \right)
        }{t,0,T}}{\mathbf{x}, \Omega}
    }{i,1,2}}{s,1,N_{s}}\\
  -
    \Sum{
        \Int{\Int{
            \eta_{s} \left( \pderiv[1]{h_{s}}{t} \right)
        }{t,0,T}}{\mathbf{x}, \Omega}
    }{s,1,N_{s}}\\
 -
    \Sum{
        \Int{\Int{
            \gamma_{s}(\mathbf{x}, t) h_{s}(\mathbf{x}, t)
        }{t, 0, T}}{\mathbf{x}, \partial\Omega}
    }{s,1,N_{s}}
  -
    \Sum{
        \Int{
            \mu_{s}(\mathbf{x}) h_{s}(\mathbf{x}, 0)
        }{t, 0, T}
    }{s,1,N_{s}}.
\end{multline}
The equation above is a derivative with relation to all possible $p_s$, for
every shot index ($p_1, p_2, p_3$), whereas we only need the derivative for
a specific one for the sought after adjoint state equation. Hence, it
is strategic here to use the Kronecker delta to enforce that inside the
definition of
% TODO put a reference to the Kronecker delta here
the test function $h_{s}$ like so:
\begin{equation}
    h_{s}(\mathbf{x}, t) = \delta_{s, \overline{s}} \, \psi(\mathbf{x}, t),
\end{equation}
which when substituted back into Eq.~\ref{eq:dLdp} results into
\begin{multline}\label{eq:dLdp1}
    \pderiv{\mathcal{L}\left(c, \mathbf{R}, p_{s}, \mathbf{v}_{s}, \mathbf{\xi}_{s}, \eta_{s}, \gamma_{s}, \mu_{s}, \mathbf{\zeta}_{s}, \mathbf{\lambda_{s}} \right)}{p_{\overline{s}}}\, \left( \delta_{s, \overline{s}} \, \psi(\mathbf{x}, t) \right) = \\
    \pderiv{J(c, \mathbf{R}, p_{s}, \mathbf{v}_{s})}{p_{\overline{s}}}\, \left( \delta_{s, \overline{s}} \, \psi(\mathbf{x}, t) \right)\\
    - 
    \Sum{
        \Int{\Int{
            \xi_{\overline{s},i} \left( \pderiv[1]{\psi}{x_{i}} \right)
        }{t,0,T}}{\mathbf{x}, \Omega}
    }{i,1,2}
  -
        \Int{\Int{
            \eta_{\overline{s}} \left( \pderiv[1]{\psi}{t} \right)
        }{t,0,T}}{\mathbf{x}, \Omega}\\
 -
        \Int{\Int{
            \gamma_{\overline{s}}(\mathbf{x}, t) \psi(\mathbf{x}, t)
        }{t, 0, T}}{\mathbf{x}, \partial\Omega}
  -
        \Int{
            \mu_{\overline{s}}(\mathbf{x}) \psi(\mathbf{x}, 0)
        }{t, 0, T},
\end{multline}
where now the derivative is in relation to a specific shot index $\overline{s}$.
The next step is to apply integration by parts in the relevant integrals to
move every differential operator away from the test function $\psi$:
\begin{multline}\label{eq:dLdp2}
    \pderiv{\mathcal{L}\left(c, \mathbf{R}, p_{s}, \mathbf{v}_{s}, \mathbf{\xi}_{s}, \eta_{s}, \gamma_{s}, \mu_{s}, \mathbf{\zeta}_{s}, \mathbf{\lambda_{s}} \right)}{p_{\overline{s}}}\, \left( \delta_{s, \overline{s}} \, \psi(\mathbf{x}, t) \right) = \\
    \pderiv{J(c, \mathbf{R}, p_{s}, \mathbf{v}_{s})}{p_{\overline{s}}}\, \left( \delta_{s, \overline{s}} \, \psi(\mathbf{x}, t) \right)\\
    - 
    \Sum{
        \Int{\Int{
            \xi_{\overline{s},i} \; \psi
        }{t,0,T}}{\mathbf{x}, \partial\Omega}
    }{i,1,2}
    + 
    \Sum{
        \Int{\Int{
            \left( \pderiv[1]{\xi_{\overline{s},i}}{x_{i}} \right) \; \psi
        }{t,0,T}}{\mathbf{x}, \Omega}
    }{i,1,2}\\
 -
        \Int{
                \eta_{\overline{s}}(\mathbf{x}, T) \psi(\mathbf{x}, T)
        }{\mathbf{x}, \Omega}\\
 +
        \Int{
                \eta_{\overline{s}}(\mathbf{x}, 0) \psi(\mathbf{x}, 0)
        }{\mathbf{x}, \Omega}
 +
        \Int{\Int{
             \left( \pderiv[1]{\eta_{\overline{s}}}{t} \right) \psi
        }{t,0,T}}{\mathbf{x}, \Omega}\\
 -
        \Int{\Int{
            \gamma_{\overline{s}}(\mathbf{x}, t) \psi(\mathbf{x}, t)
        }{t, 0, T}}{\mathbf{x}, \partial\Omega}
 -
        \Int{
            \mu_{\overline{s}}(\mathbf{x}) \psi(\mathbf{x}, 0)
        }{t, 0, T},
\end{multline}
and here we can define the adjoint variables for the boundary conditions to be
$\gamma_{\overline{s}}(\mathbf{x}, t) = - \Sum{\xi_{\overline{s},i}}{i,1,2}(\mathbf{x}, t)$,
and $\mu_{\overline{s}}(\mathbf{x}) = \eta_{\overline{s}}(\mathbf{x}, 0)$. That
way, the boundary and initial terms cancel if we also set 
$\eta_{s}(\mathbf{x}, T) = 0$.
\begin{multline}\label{eq:dLdp3}
    \pderiv{\mathcal{L}\left(c, \mathbf{R}, p_{s}, \mathbf{v}_{s}, \mathbf{\xi}_{s}, \eta_{s}, \gamma_{s}, \mu_{s}, \mathbf{\zeta}_{s}, \mathbf{\lambda_{s}} \right)}{p_{\overline{s}}}\, \left( \delta_{s, \overline{s}} \, \psi(\mathbf{x}, t) \right) = \\
    \pderiv{J(c, \mathbf{R}, p_{s}, \mathbf{v}_{s})}{p_{\overline{s}}}\, \left( \delta_{s, \overline{s}} \, \psi(\mathbf{x}, t) \right)\\
    + 
    \Sum{
        \Int{\Int{
            \left( \pderiv[1]{\xi_{\overline{s},i}}{x_{i}} \right) \; \psi
        }{t,0,T}}{\mathbf{x}, \Omega}
    }{i,1,2}
 +
        \Int{\Int{
             \left( \pderiv[1]{\eta_{\overline{s}}}{t} \right) \psi
        }{t,0,T}}{\mathbf{x}, \Omega}\\
\end{multline}
which leave us with the only with the derivative of the cost function. This
derivative is calculated as follows:
\begin{multline}\label{eq:costDeriv}
    \pderiv{J(c, \mathbf{R}, p_{s}, \mathbf{v}_{s})}{p_{\overline{s}}}\, \left( \delta_{s, \overline{s}} \, \psi(\mathbf{x}, t) \right)=\\
    \Bigg[ \Dt{\epsilon} \frac{1}{\beta} \Sum{ \Sum{
            \Int{
                \Big( S_{s,r} \left( p_{s}(\mathbf{x}, t) + \epsilon \, \delta_{s, \overline{s}} \, \psi(\mathbf{x}, t) \right) - d_{s,r}(t) \Big)^2
            }{t, 0, T} 
        }{r, 1, N_{r}} }{s, 1, N_{s}} \Bigg]_{\epsilon=0},
\end{multline}
with $\beta=\frac{1}{2 N_{s} N_{r} T}$. Given this is a derivative to $\epsilon \in \mathbb{R}$,
it can go inside the sums and integrals, leading to:
\begin{multline}\label{eq:costDeriv1}
    \pderiv{J(c, \mathbf{R}, p_{s}, \mathbf{v}_{s})}{p_{\overline{s}}}\, \left( \delta_{s, \overline{s}} \, \psi(\mathbf{x}, t) \right)=\\
    \frac{1}{\beta} \Sum{
            \Int{
                2 \Big( S_{\overline{s},r} \left( p_{\overline{s}}(\mathbf{x}, t) \right) - d_{\overline{s},r}(t) \Big) S_{\overline{s},r}  \, \psi(\mathbf{x}, t)
            }{t, 0, T} 
        }{r, 1, N_{r}}.
\end{multline}
The objective here is to put Eq.~\ref{eq:costDeriv1} with the same space-time 
integrals of the other terms in Eq.~\ref{eq:dLdp3}, which will require playing
around with the sampling operator definition:
\begin{multline}\label{eq:costDeriv2}
    \pderiv{J(c, \mathbf{R}, p_{s}, \mathbf{v}_{s})}{p_{\overline{s}}}\, \left( \delta_{s, \overline{s}} \, \psi(\mathbf{x}, t) \right)\\=
    \frac{1}{\beta} \Sum{
            \Int{
                2 \Big( S_{\overline{s},r} \left( p_{\overline{s}}(\mathbf{x'}, t) \right) - d_{\overline{s},r}(t) \Big) \Int{\psi(\mathbf{x}, t) \delta(\mathbf{x} - \mathbf{x}_{s,r})}{\mathbf{x}, \Omega}
            }{t, 0, T} 
        }{r, 1, N_{r}}
                \\ =
    \frac{2}{\beta}  \Sum{
             \Int{
                 \Int{\Big( S_{\overline{s},r} \left( p_{\overline{s}}(\mathbf{x'}, t) \right) - d_{\overline{s},r}(t) \Big) \delta(\mathbf{x} - \mathbf{x}_{s,r})\, \psi(\mathbf{x}, t)}{\mathbf{x}, \Omega}
            }{t, 0, T} 
        }{r, 1, N_{r}}
                \\ =
    \frac{2}{\beta} \Sum{
             \Int{\Int{
                 S_{\overline{s},r}^{\dagger} \Big( S_{\overline{s},r} \left( p_{\overline{s}}(\mathbf{x'}, t) \right) - d_{\overline{s},r}(t) \Big) \, \psi(\mathbf{x}, t)
             }{\mathbf{x}, \Omega}}{t, 0, T} 
        }{r, 1, N_{r}},
\end{multline}
where $S_{\overline{s},r}^{\dagger}$ is the adjoint of the sampling operator,
which just means placing a value at the position of the receiver.
Substituting the equation above (Eq.~\ref{eq:costDeriv1}) into the expression
of the derivative of the lagrangian with relation to the pressure (Eq.~\ref{eq:dLdp3})
while equaling to zero to satisfy the adjoint state method conditions
results into:
\begin{multline}\label{eq:dLdp4}
    \pderiv{\mathcal{L}\left(c, \mathbf{R}, p_{s}, \mathbf{v}_{s}, \mathbf{\xi}_{s}, \eta_{s}, \gamma_{s}, \mu_{s}, \mathbf{\zeta}_{s}, \mathbf{\lambda_{s}} \right)}{p_{\overline{s}}}\, \left( \delta_{s, \overline{s}} \, \psi(\mathbf{x}, t) \right) = \\
    \frac{2}{\beta} \Sum{
             \Int{\Int{
                 S_{\overline{s},r}^{\dagger} \Big( S_{\overline{s},r} \left( p_{\overline{s}}(\mathbf{x'}, t) \right) - d_{\overline{s},r}(t) \Big) \, \psi(\mathbf{x}, t)
             }{\mathbf{x}, \Omega}}{t, 0, T} 
        }{r, 1, N_{r}}\\
+ 
    \Sum{
        \Int{\Int{
            \left( \pderiv[1]{\xi_{\overline{s},i}}{x_{i}} \right) \; \psi
        }{t,0,T}}{\mathbf{x}, \Omega}
    }{i,1,2}
+
        \Int{\Int{
             \left( \pderiv[1]{\eta_{\overline{s}}}{t} \right) \psi
        }{t,0,T}}{\mathbf{x}, \Omega} = 0,
\end{multline}
which is a weak/integral formulation of the following PDE:
\Style{DDisplayFunc=inset,DShorten=false}
\begin{multline}\label{eq:dLdp5}
             \pderiv[1]{\eta_{s}}{t}
+
    \Sum{
             \pderiv[1]{\xi_{s,i}}{x_{i}}
    }{i,1,2}
+ 
    \frac{2}{\beta} \Sum{
                 S_{s,r}^{\dagger} \Big( S_{s,r} \left( p_{s}(\mathbf{x}, t) \right) - d_{s,r}(t) \Big)
        }{r, 1, N_{r}}=0
\end{multline}
where the $\overline{s}$ were replaced with $s$ to simplify notation, given
now it is clear that the equation is the same for any shot index.
Calculating the derivative of the lagrangian in Eq.~\ref{eq:lagrangian}
with relation to the particle momentum $v_{s,i}$, for a specific shot index,
will result into the other adjoint state equation:
\Style{DDisplayFunc=outset,DShorten=false}
\begin{equation}\label{eq:dLdv}
    \pderiv{\xi_{s,i}}{t} +
    \pderiv{\left( c^2 \eta \right)}{x_{i}} + 
     \eta \left(2 c^2 R_{i} - c \pderiv{c}{x_{i}} \right) = 0,
\end{equation}
which together with Eq.~\ref{eq:dLdp5} compose the adjoint system of equations
to be solved for $\xi_{s,i}$ and $\eta_{s}$.

The last step is to calculate the derivative of the lagrangian
(Eq.~\ref{eq:lagrangian}) with relation to the reflectivity, to a specific
component $\overline{i}$, in a similar fashion to the other derivatives:
\begin{multline}\label{eq:dLdR}
    \pderiv{\mathcal{L}\left(c, \mathbf{R}, p_{s}, \mathbf{v}_{s}, \mathbf{\xi}_{s}, \eta_{s}, \gamma_{s}, \mu_{s}, \mathbf{\zeta}_{s}, \mathbf{\lambda_{s}} \right)}{R_{\overline{i}}} \left( \delta_{i,\overline{i}} \, \phi \right) \\=
    2 \Int{\Int{
        c^2(\mathbf{x}) \, \eta(\mathbf{x}, t) \, v_{s, \overline{i}}(\mathbf{x}, t) \,\phi(\mathbf{x})
    }{\mathbf{x}, \Omega}}{t, 0, T} \\=
    \left<\
    2 \Int{
        c^2(\mathbf{x}) \, \eta(\mathbf{x}, t) \, v_{s, \overline{i}}(\mathbf{x}, t) 
    }{t, 0, T} 
    ,\;
    \phi(\mathbf{x})
    \right>_{\Omega},
\end{multline}
where $\phi(\mathbf{x})$ is just a test function. The integral over space in
this directional functional derivative is a inner product of the functional
derivative kernel (Fréchet, non-directional) with the test function, which collapses the space
variable and makes the resulting test derivative be a single scalar for the whole
space $\Omega$. When discretized, this non-directional derivative becomes
a jacobian, while the directional version becomes the dot product of the
jacobian with a test vector. Strategically here, a simple way to obtain
the derivative one can simply choose the test function to be a Dirac delta
at another point in space, and use the its sampling property:
\begin{equation}
    \phi(\mathbf{x}) = \delta(\mathbf{x} - \mathbf{x}'),
\end{equation}
which when applied to the Eq.~\ref{eq:dLdR} leads to the sought after
derivative of the cost function:
\begin{equation}\label{eq:gradient}
    \CoolD[1]{J}{R_{i}} = 
2 \Int{
        c^2(\mathbf{x}) \, \eta(\mathbf{x}, t) \, v_{s,i}(\mathbf{x}, t) 
    }{t, 0, T},
\end{equation}
where again it is clear that the expression for the derivative is the
same for different components ($i$), hence the equation above has
$i$ instead of $\overline{i}$. The equation above (Eq.~\ref{eq:gradient}),
when discretized is the gradient of the cost function with relation to the
vector-reflectivity.


%%%%%%%%%%%%  Supplementary Figures  %%%%%%%%%%%%
%\clearpage

%%%%%%%%%%%%%%%%   End   %%%%%%%%%%%%%%%%
%\end{multicols}  % Method B for two-column formatting (doesn't play well with line numbers), comment out if using method A
\end{document}
